\section*{अध्याय \ref{ch:b2:plant-meristems} --- पौधों का कायिक विकास: विभज्योतक और वृद्धि}

\begin{solution}{exo:b2:plant-meristems:1}
\omterm{def:b2:plant-meristems:meristem}{विभज्योतक} अविभेदित, विभाजित होती कोशिकाओं की ऐसी टिकाऊ आबादी है
जिससे उस पौधे के ऊतक निकलते हैं। \omterm{def:b2:plant-meristems:meristem}{प्ररोह शीर्षस्थ विभज्योतक}:
पत्तियाँ, तना, \omterm{def:b2:angiosperm-reproduction:flower}{पुष्प}; \omterm{def:b2:plant-meristems:meristem}{मूल शीर्षस्थ विभज्योतक}: जड़; \omterm{def:b2:plant-meristems:wood}{संवहनी
कैंबियम}: द्वितीयक जाइलम (\omterm{def:b2:plant-meristems:wood}{काष्ठ}) और फ्लोएम; कॉर्क \omterm{def:b2:plant-meristems:meristem}{कैंबियम}:
\omterm{def:b2:plant-meristems:wood}{छाल} का \omterm{def:b2:plant-meristems:wood}{परित्वक्} (कॉर्क)।
\end{solution}

\begin{solution}{exo:b2:plant-meristems:2}
\omterm{prop:b2:plant-meristems:ram}{मूल टोप} (रक्षा, चिकनाई, गुरुत्व-संवेदन, कोशिकाएँ झड़ती हुईं);
\omterm{prop:b2:plant-meristems:ram}{निष्क्रिय केंद्र} सहित \omterm{def:b2:plant-meristems:meristem}{विभज्योतक} (विभाजन); \omterm{prop:b2:plant-meristems:ram}{दीर्घीकरण क्षेत्र}
(कोशिकाएँ दस से बीस गुना लंबी होती हैं); \omterm{def:b2:cell-differentiation:differentiation}{विभेदन} क्षेत्र (मूलरोम,
परिपक्व जाइलम तथा फ्लोएम)। पार्श्व जड़ें परिपक्व जड़ के
परिरंभ से, जाइलम ध्रुवों के सामने, उठती हैं।
\end{solution}

\begin{solution}{exo:b2:plant-meristems:3}
सेम के पौधों का शीर्ष काट दिया गया: पार्श्व कलिकाएँ निकल आईं। दूसरे
समूह में कटे ठूँठ को ऑक्सिनयुक्त अगार का एक खंड दिया गया:
कलिकाएँ प्रसुप्त रहीं; और नियंत्रण समूह पर सादे अगार का खंड: कलिकाएँ
निकल आईं। निष्कर्ष: शीर्ष की ऑक्सिन ही पार्श्व कलिकाओं को संदमित करती
है --- यानी \omterm{prop:b2:plant-meristems:apical}{शीर्ष प्रभाविता} हॉर्मोनी है।
\end{solution}

\begin{solution}{exo:b2:plant-meristems:4}
साठ वर्ष की आयु; \omterm{def:b2:plant-meristems:wood}{रस-काष्ठ} अंतिम पंद्रह वलय हैं; और जीवित हैं
\omterm{def:b2:plant-meristems:meristem}{विभज्योतक}, फ्लोएम, \omterm{def:b2:plant-meristems:wood}{रस-काष्ठ} की रश्मि तथा मृदूतक कोशिकाएँ,
और कॉर्क \omterm{def:b2:plant-meristems:meristem}{विभज्योतक} --- यानी किसी मृत गूदे के चारों ओर एक पतला खोल।
\end{solution}

\begin{solution}{exo:b2:plant-meristems:5}
$0.1\times(P - 0.3)$: \qty{0.01}{h^{-1}} (दुगुना होने में $\ln 2/0.01 =
\qty{69}{h}$), \qty{0.03}{h^{-1}} (\qty{23}{h}), \qty{0.05}{h^{-1}}
(\qty{14}{h})। \qty{0.25}{MPa} पर स्फीति पराभव देहली से
नीचे है: वृद्धि रुक जाती है, किसी भी मुरझान से पहले।
\end{solution}

\begin{solution}{exo:b2:plant-meristems:6}
पत्तियाँ 1--8: 0, 137.5, 275, 52.5, 190, 327.5, 105, \qty{242.5}{\degree}।
क्रम में: 0, 52.5, 105, 137.5, 190, 242.5, 275, 327.5 --- यानी सबसे छोटी
ख़ाली जगह \qty{32.5}{\degree}। हर पत्ती पिछली पत्तियों की किसी ख़ाली जगह
में बैठती है, इसलिए वे आठों एक-दूसरे पर जितनी कम हो सके उतनी कम छाया
डालती हैं और वह मुकुट प्रकाश की चक्रिका समान रूप से भर देता है।
\end{solution}

\begin{solution}{exo:b2:plant-meristems:7}
$2\pi r h\Delta r$: \qty{5}{cm} पर $2\pi\times 0.05\times 12\times
0.003 = \qty{0.011}{m^3}$, \qty{5.7}{kg}, \qty{2.8}{kg} कार्बन; और
\qty{40}{cm} पर \qty{0.090}{m^3}, \qty{45}{kg}, \qty{23}{kg} कार्बन
--- यानी उसी वलय-चौड़ाई से आठ गुना अधिक।
\end{solution}

\begin{solution}{exo:b2:plant-meristems:8}
\emph{clv3}: कोई ब्रेक नहीं --- विशाल \omterm{def:b2:plant-meristems:meristem}{विभज्योतक}, अतिरिक्त पत्तियाँ तथा
पुष्प-अंग, फ़ैसिएटेड तने। \emph{wus}: कोई त्वरक नहीं --- वह \omterm{def:b2:plant-meristems:meristem}{विभज्योतक}
कुछ पत्तियों के बाद चुक जाता है और बार-बार फिर बसाया जाता है। दोहरा
उत्परिवर्ती: \emph{wus} का \omterm{def:b2:meiosis-heredity:vocabulary}{लक्षण-प्ररूप}, इसलिए WUS CLV3 के नीचे काम करता है:
CLV3 का काम WUS को दबाना है, और WUS ही न हो तो दबाने को कुछ नहीं बचता।
\end{solution}

\begin{solution}{exo:b2:plant-meristems:9}
अभिगमन: \qty{20}{h}। विसरण: $(0.2)^{2}/(2\times 10^{-9}) =
\qty{2e7}{s}$, यानी लगभग 230 दिन। ध्रुवीय अभिगमन किसी हॉर्मोनी
संकेत को दिन भर में पूरे पौधे की लंबाई पर काम का बना देता है, और उसे
एक दिशा भी देता है।
\end{solution}

\begin{solution}{exo:b2:plant-meristems:10}
\qty{200}{\micro m} की कोशिकाओं वाले परिपक्व ऊतक का प्रतिदिन \qty{3}{cm}:
यानी प्रति पंक्ति प्रतिदिन \num{150} कोशिकाएँ। इसलिए सौ विभाजित होती
कोशिकाओं में से हर एक को प्रतिदिन 1.5 बार विभाजित होना पड़ता है: यानी लगभग \qty{16}{h} का चक्र।
\end{solution}

\begin{solution}{exo:b2:plant-meristems:11}
जिबरेलिन पर्वसंधि-अंतराल लंबे करता है; और जो पौधा उत्तर ही न दे सके उसके
पर्वसंधि-अंतराल छोटे और तना छोटा होता है। भारी नाइट्रोजन खाद किसी
लंबे गेहूँ को और लंबा तथा भारी-बाली वाला बना देती है जब तक वह गिर
(लॉज) न जाए और सड़ न जाए; जबकि छोटा कड़ा तना खड़ा रहता है और अतिरिक्त
कार्बन पुआल के बजाय दाने में डालता है, इसलिए कटाई-सूचकांक चढ़ जाता है।
जंगल में कोई छोटा पौधा अपने पड़ोसियों से ऊपर से ढक और छा लिया जाता है
और हार जाता है।
\end{solution}

\begin{solution}{exo:b2:plant-meristems:12}
\omterm{def:b2:plant-meristems:meristem}{विभज्योतक} किसी पौधे को वहीं और तभी अंग जोड़ने देते हैं जहाँ और जब
परिस्थितियाँ अनुकूल हों --- पत्तियाँ प्रकाश की ओर, जड़ें जल की ओर --- और
जो न फलें उन्हें छोड़ देने देते हैं; हर कक्षस्थ कलिका एक अतिरिक्त
वृद्धि-बिंदु है, इसलिए चराई, आग और छँटाई पुनर्वृद्धि से झेल ली जाती हैं;
और चूँकि वे \omterm{def:b2:plant-meristems:meristem}{विभज्योतक} सदा भ्रूणीय रहते हैं, किसी \omterm{def:b2:asexual-reproduction:clone}{जेनेट} की कोई
नियत आयु नहीं होती। क़ीमत: कोई पौधा न चल सकता है न जो गढ़ चुका उसे फिर से
सजा सकता है --- एक बार रखा गया अंग वहीं रहता है, और संचयन से गढ़ी गई कोई
काया अपना मृत अतीत (अंतःकाष्ठ) साथ ढोती है।
\end{solution}

\begin{solution}{pb:b2:plant-meristems:1}
\textbf{1.} प्रति शीर्ष $150/3 = 50$ पत्तियाँ; और उस वृक्ष के लिए \num{10000}।
\textbf{2.} 0, 137.5, 275, 52.5 और \qty{190}{\degree}।
\textbf{3.} $8\times 137.5 = 1100 = 3\times 360 + 20$; $13\times 137.5
= 1787.5 = 5\times 360 - 12.5$; $21\times 137.5 = 2887.5 = 8\times 360
+ 7.5$: यानी 22वीं पत्ती पहली ऐसी है जो पहली पत्ती के
\qty{10}{\degree} के भीतर पड़ती है।
\textbf{4.} हर नई पत्ती सबसे चौड़ी ख़ाली जगह में बैठती है, इसलिए पत्तियाँ
एक-दूसरे पर सबसे कम छाया डालती हैं; और वह ऑक्सिन ही है जिसका ध्रुवीय
अभिगमन (PIN वाहक) हर प्राथमिकांग के परिवेश को ख़ाली कर देता है, और
यही अंतराल तय करता है।
\textbf{5.} $\num{10000}\times 80\,\text{cm}^{2} = \qty{80}{m^2}$: यानी
पूरा मुकुट इसी ऋतु की पत्तियाँ हैं।
\textbf{6.} पूरा; और पहली पत्तियाँ पिछले वर्ष \omterm{def:b2:plant-meristems:wood}{काष्ठ} तथा जड़ों में सँजोए
मंड से गढ़ी जाती हैं।
\textbf{7.} तिगुना शीर्ष प्राथमिकांग तेज़ बनाता --- प्रति
प्लास्टोक्रॉन दो या तीन --- इसलिए पत्तियाँ और पर्ण-क्षेत्रफल तिगुने हो
सकते, पर वे प्ररोह फ़ैसिएटेड होते, उनके तने फ़ीते जैसे और पर्णविन्यास
अव्यवस्थित होता, और उस अतिरिक्त पत्ते का अधिकांश स्वयं पर छाया डालता।
\textbf{8.} दिन: $0.05\times(0.55 - 0.35) = \qty{0.01}{h^{-1}}$; रात:
$0.05\times 0.40 = \qty{0.02}{h^{-1}}$।
\textbf{9.} दिन: $2e^{0.12} = \qty{2.25}{cm}$, यानी \qty{0.25}{cm} जुड़े;
रात: $2.25e^{0.24} = \qty{2.87}{cm}$, यानी \qty{0.61}{cm} जुड़े।
\textbf{10.} \qty{96}{h} पर माध्य दर \qty{0.015}{h^{-1}}: $2e^{1.44}
= \qty{8.4}{cm}$।
\textbf{11.} दिन: $0.05\times 0.30 = \qty{0.015}{h^{-1}}$; रात:
$0.05\times 0.50 = \qty{0.025}{h^{-1}}$।
\textbf{12.} $0.05\times 0.05 = \qty{0.0025}{h^{-1}}$, यानी सामान्य
दिन-दर का चौथाई। परासरणी समायोजन उस कोशिका का विलेय-अंश चढ़ा देता है
ताकि उसी जल विभव पर स्फीति फिर उस देहली से ऊपर चढ़ जाए।
\textbf{13.} दिन में वाष्पोत्सर्जन उस पत्ती का जल विभव और उसके साथ
स्फीति गिरा देता है, इसलिए $P - Y$ छोटा रहता है; और रात में स्फीति
बहाल हो जाती है। वृद्धि घंटे-दर-घंटे जल से सीमित है, शर्करा से नहीं।
\textbf{14.} $2\pi\times 0.10\times 8\times 0.005 = \qty{0.025}{m^3}$;
\qty{12.6}{kg} शुष्क।
\textbf{15.} \qty{6.3}{kg} कार्बन।
\textbf{16.} $80\times 4\times 150 = \qty{48}{kg}$ कार्बन।
\textbf{17.} $6.3/48 = \qty{13}{\%}$। और सिंक: स्वयं पत्तियाँ,
जड़ें, शाखाएँ तथा टहनियाँ, हर जीवित कोशिका का श्वसन (जो स्थिर किए गए का
लगभग आधा है), भंडार, \omterm{def:b2:angiosperm-reproduction:flower}{पुष्प} तथा
\omterm{def:b2:angiosperm-reproduction:seed}{बीज}।
\textbf{18.} त्रिज्या \qty{0.20}{m}: $2\pi\times 0.20\times 8\times
0.005 = \qty{0.050}{m^3}$, यानी इस वर्ष का दुगुना --- क्योंकि जिस परिधि
पर \omterm{def:b2:plant-meristems:meristem}{विभज्योतक} काम करता है वह दुगुनी हो चुकी है (और वह वृक्ष ऊँचा भी होगा)।
\textbf{19.} कोई शुष्क या ठंडा वर्ष; या कीटों, पाले या आग से पर्णलोप।
कोई जलवायु-कारण उस क्षेत्र के हर वृक्ष में वही वलय सँकरे करता है और
मौसम-अभिलेखों से मेल खाता है; जबकि क्षति एक ही वृक्ष या झुंड को छूती है,
और पाला या आग उस वलय में शारीरिक निशान छोड़ जाते हैं।
\textbf{20.} हर कटान के सबसे पास की कलिकाएँ दिनों में कई प्ररोहों में
निकल आती हैं; और महीने भर में वह वृक्ष अधिक पर छोटे प्ररोहों के साथ
झाड़ीदार तथा सघन हो जाता है।
\textbf{21.} कटान पर ऑक्सिन शीर्ष के संकेत की जगह ले लेती है: कलिकाएँ
प्रसुप्त रहती हैं और कोई शाखन नहीं होता।
\textbf{22.} हर कटाई शीर्ष हटा देती है और नीचे की कलिकाएँ मुक्त कर देती
है; वर्ष में छह बार, और हर मुक्ति छोटे प्ररोहों की एक परत जोड़ देती है,
और वह सतह उन्हीं से भर जाती है।
\textbf{23.} ठूँठ की प्रसुप्त कक्षस्थ कलिकाओं से और \omterm{def:b2:plant-meristems:meristem}{विभज्योतक} की बनाई
अपस्थानिक कलिकाओं से; और हर कलिका एक पूरा \omterm{def:b2:plant-meristems:meristem}{विभज्योतक} है जो
प्ररोह-तंत्र फिर गढ़ सकता है, जबकि किसी प्राणी का सिर एक ही, न बदला जा
सकने वाला \omterm{def:b2:vertebrate-development:induction}{संघटक} है जिसकी कोई आरक्षित प्रति नहीं।
\textbf{24.} कलिकाओं तक अधिक साइटोकाइनिन पहुँचना उनके निकलने को बढ़ावा
देता है --- यानी छँटे हुए पौधे को पानी देने से वह अधिक शाखित होता है।
\textbf{25.} प्रति शीर्ष 50 पत्तियाँ, प्रति वृक्ष \num{10000};
पर्वसंधि-अंतराल की दरें दिन में \qty{0.01}{h^{-1}} और रात में \qty{0.02}{h^{-1}};
और वर्ष की \omterm{def:b2:plant-meristems:wood}{काष्ठ} \qty{0.025}{m^3}, \qty{12.6}{kg} शुष्क, \qty{6.3}{kg} कार्बन।
\end{solution}
